% 图 2：一般线性离散动态系统的矩阵框图
\documentclass[tikz,border=4pt]{standalone}
\input{common}
\input{kalman-style}
\begin{document}
\begin{tikzpicture}[kalman]
  \coordinate (in) at (0,0);
  \node[block] (D)  at (1.6,0) {$\bm{\Delta}(t)$};
  \node[sum]   (s)  [right=10mm of D] {$\Sigma$};
  \node[tri]   (z)  [right=18mm of s] {\TXTdelay};
  \coordinate  (br) at ($(z.east)+(22mm,0)$);
  \node[block] (M)  [right=10mm of br] {$\bm{M}(t)$};
  \coordinate  (out) at ($(M.east)+(12mm,0)$);
  \sumsigns{s}

  \draw[sig] (in) node[left,lbl]{$\bm{u}(t)$} -- (D);
  \draw[sig] (D) -- (s);
  \draw[sig] (s) -- node[above,lbl]{$\bm{x}(t+1)$} (z.west);
  \draw[edge] (z.east) -- node[above,lbl]{$\bm{x}(t)$} (br);
  \draw[sig] (br) -- (M);
  \draw[sig] (M) -- (out) node[right,lbl]{$\bm{y}(t)$};
  \node[dot] at (br) {};

  % 反馈：x -> Phi(t+1; t) -> Sigma
  \node[block] (Phi) at ($(z)+(0,-15mm)$) {$\bm{\Phi}(t+1;t)$};
  \draw[sig] (br) |- (Phi);
  \draw[sig] (Phi) -| (s);
\end{tikzpicture}
\end{document}
